\documentclass[conference]{IEEEtran}
\IEEEoverridecommandlockouts
\usepackage{cite}
\usepackage{amsmath,amssymb}
\usepackage{graphicx}
\usepackage{booktabs}
\usepackage{xcolor}
\usepackage{circuitikz}
\usepackage[caption=false,font=footnotesize]{subfig}
\usetikzlibrary{arrows.meta,positioning,calc,shapes.geometric}
\graphicspath{{figures/}{figures/revb/}}
\newcommand{\uF}{\ensuremath{\mu}F}
\input{paper_numbers.tex}
\input{device_table.tex}

\begin{document}

\title{Building Block Design for Stacked Polyphase Bridge Converters for High-Voltage Traction Drives}

\author{\IEEEauthorblockN{Enes Ayaz}
\IEEEauthorblockA{\textit{Division of Electric Power and Energy Systems} \\
\textit{KTH Royal Institute of Technology}\\
Stockholm, Sweden}
% TODO: co-authors and e-mail addresses
}

\maketitle

\begin{abstract}
The stacked polyphase bridge (SPB) converter divides a high traction-battery voltage among series-connected three-phase submodules, so that the transistor voltage class becomes a design variable instead of a consequence of the battery voltage. This paper derives the phase-leg building block of a 1200\,V, 300\,kW SPB inverter step by step, from the specification to hardware. Because the phase current, \Irms\,A$_\mathrm{rms}$ per leg, does not depend on the number of submodules, the device time constant $R_\mathrm{DS(on)}C_\mathrm{oss,tr}$ and the on-resistance per volt of blocking capability select 100\,V GaN transistors, and the overshoot $L_\mathrm{loop}\,di/dt$ at 0.5 and 1\,nH sets 16 submodules of 75\,V. Loss, junction temperature and peak voltage give four EPC2361 per switch, and a current-sharing sensitivity analysis gives one dual-channel isolated driver per pair of transistors in two mirrored cells. The dc link of a three-phase level is sized for 3\,\% voltage ripple with the measured dc-bias derating of class-II MLCCs, which selects twenty 22\,\uF{} 2220 capacitors per leg. The resulting four-cell leg board is built in two assemblies, one with a coaxial shunt in every cell for double-pulse testing and one without shunt for continuous operation; three legs share a laminated PCB busbar. Q3D extraction of the routed boards gives a commutation loop of \LoopLoad\,nH per cell without and \LoopDpt\,nH with the shunt, equal in all four cells within \LoopDptSpread\,\%, and \BusLab\,nH between neighbouring legs on the busbar.
\end{abstract}

\begin{IEEEkeywords}
Gallium nitride, dc-link capacitor, paralleling, parasitic inductance, PCB busbar, stacked polyphase bridge, traction inverter.
\end{IEEEkeywords}

%=========================================================================
\section{Introduction}
%=========================================================================
Raising the battery voltage of electric vehicles to 800\,V and beyond reduces the current and the charging time, but it forces a two-level traction inverter to 1200\,V transistors, steep voltage steps at the machine terminals and bulky dc-link capacitors. The stacked polyphase bridge (SPB) converter~\cite{norrga2013,jin2014} offers an alternative: $N$ three-phase bridges are connected in series on the dc side, each with its own dc link and each feeding one three-phase winding set of a multiphase machine. Every submodule sees only $V_\mathrm{bus}/N$, so low-voltage transistors with a far better figure of merit can be used, the submodule capacitors shrink with a higher switching frequency~\cite{jin2017}, the voltage step at the windings is reduced~\cite{bringezu2023}, and a faulty submodule can be bypassed~\cite{jin2017}. The dc-link stability~\cite{nikouie2017tpel}, control~\cite{jin2014,bringezu2023,verkroost2024}, machine design~\cite{zhang2014,bringezu2020} and an integrated prototype~\cite{nikouie2017} have been reported, and a comparative study found a GaN-based SPB competitive with a three-level flying-capacitor inverter for integrated motor drives~\cite{rohner2024}.

The building block of an SPB is nevertheless a demanding power stage in its own right. At 1200\,V and 300\,kW each phase leg carries more than 200\,A$_\mathrm{rms}$ at only 75\,V, so several GaN transistors must be paralleled, and paralleled GaN half bridges are known to fail by drain--source overshoot~\cite{brothers2019}, by uneven dynamic current sharing in asymmetric loops~\cite{lu2017,tran2025} and by gate-voltage disturbances~\cite{tran2025,musumeci2022}. Low-voltage paralleled-GaN inverters have been demonstrated with symmetric layouts~\cite{wattenberg2023,cooke2026} and hybrid packaging~\cite{tran2025}, and efficiency--density co-optimization has been formalized for a 100\,kW GaN inverter~\cite{zou2025}. What is missing for the SPB is a procedure in which the submodule voltage, the transistor, the number of paralleled devices, the gate drive and the dc link follow from each other and from the inverter specification, instead of being fixed in advance, and which is carried through to a board that can be built and tested.

This paper presents such a procedure (Fig.~\ref{fig:flow}) for a 1200\,V, 300\,kW SPB and the resulting hardware. Its contributions are:
\begin{enumerate}
\item a voltage-class selection in which the minimum semiconductor loss scales with the device time constant and the conduction loss with the on-resistance per volt of blocking capability, both of which favor 100\,V GaN transistors;
\item a submodule-voltage and device-count derivation that treats the commutation-loop inductance (0.5 and 1\,nH) and the thermal path of the test hardware explicitly;
\item a gate-drive concept with one dual-channel isolated driver per pair of transistors, justified by a current-sharing sensitivity analysis, and a dc-link design for a three-phase level that includes the dc-bias derating of class-II MLCCs;
\item a four-cell leg board in two assemblies (with and without current shunts), its verification by Q3D extraction of the routed copper, and a laminated PCB busbar that joins three legs to one level.
\end{enumerate}

%=========================================================================
\section{SPB Inverter and Design Procedure}
%=========================================================================
\begin{figure*}[t]
\centering
\resizebox{\textwidth}{!}{\input{figures/fig_spb_b.tex}}
\caption{(a) SPB inverter: $N$ three-phase submodules (SM) in series on the dc side feed the $N$ winding sets of a multiphase machine; each submodule operates at $V_m=V_\mathrm{bus}/N$. (b) Building block designed in this paper, one phase leg of a submodule: four commutation cells, each with a high-side (QH) and a low-side (QL) EPC2361 and its bypass capacitors $C_\mathrm{byp}$, a common bank $C_\mathrm{bank}$, one dc terminal pair and one ac output; one dual-channel driver serves each pair of cells through equal gate stars. The coaxial shunt (dashed) is fitted on the double-pulse-test board only. (c) Gate and power circuit of one transistor with the common-source inductance $L_{CS}$ up to the Kelvin connection K.}
\label{fig:spb}
\end{figure*}

If the bus voltage $V_\mathrm{bus}$ is divided into $N$ submodules of $V_m=V_\mathrm{bus}/N$ that share the power equally (Fig.~\ref{fig:spb}a), one phase leg delivers at modulation index $M$ and power factor $\cos\varphi$
\begin{equation}
P_\mathrm{leg}=\frac{M\,V_m}{2\sqrt2}I_\mathrm{ph}\cos\varphi,\quad P=3NP_\mathrm{leg}=\frac{3M\,V_\mathrm{bus}}{2\sqrt2}I_\mathrm{ph}\cos\varphi.
\label{eq:pleg}
\end{equation}
The phase current is therefore set by the inverter power and the bus voltage alone and is the same for every choice of $N$. With space-vector modulation ($M$\,=\,1.15) and $\cos\varphi$\,=\,0.9 of a permanent-magnet machine, 300\,kW at 1200\,V requires \Irms\,A$_\mathrm{rms}$ (\Ipk\,A peak) per leg. The phase current is the fixed specification of the building block; its voltage, the transistor, the number of paralleled devices and the dc link remain open and are fixed in the order of Fig.~\ref{fig:flow}. Two further targets complete the specification: the peak drain--source voltage must stay below 90\,V, 10\,\% below the rating of the 100\,V class, and the switching frequency must be at least 50\,kHz to keep the dc links and the current ripple of the windings small.

\begin{figure*}[t]
\centering
\resizebox{\textwidth}{!}{\input{figures/fig_flow_b.tex}}
\caption{Design procedure of the SPB building block. Steps 1--6 fix the device-level quantities, steps 7--10 the board and the system integration, step 11 verifies the routed copper; the dashed paths return to the layout, the device count, the dc link or the switching speed if a target is missed.}
\label{fig:flow}
\end{figure*}

%=========================================================================
\section{Transistor and Submodule Voltage}\label{sec:vclass}
%=========================================================================
\subsection{Voltage class}
Every submodule carries the full phase current $I$, see~\eqref{eq:pleg}. With $n$ transistors per switch, the conduction loss of all submodules scales with $NI^2R_\mathrm{DS(on)}/n$ and the loss associated with the output charge with $Nnf_\mathrm{sw}Q_\mathrm{oss}V_m$. With $Q_\mathrm{oss}\approx C_\mathrm{oss,tr}V_m$ and the optimal number of devices $n^\ast=I\sqrt{R_\mathrm{DS(on)}/(f_\mathrm{sw}Q_\mathrm{oss}V_m)}$, their sum over all legs becomes
\begin{equation}
P_\mathrm{min}=6\,I\,V_\mathrm{bus}\sqrt{f_\mathrm{sw}\,R_\mathrm{DS(on)}C_\mathrm{oss,tr}} .
\label{eq:pmin}
\end{equation}
To first order the voltage class cancels, and candidates are ranked by the device time constant $\tau=R_\mathrm{DS(on)}C_\mathrm{oss,tr}$. Fig.~\ref{fig:devices}(a) and Table~\ref{tab:devices} compare ten off-the-shelf GaN transistors of the 100, 200, 350 and 650\,V classes, each operated at 60--75\,\% of its rating. The 100\,V parts reach the lowest $\tau$, and EPC2361 the lowest of all. Normalized to the blocking voltage (Fig.~\ref{fig:devices}b), the 100\,V class has the lowest on-resistance per volt (7.5\,\textmu$\Omega$/V for EPC2361, three to thirty times below the 200--650\,V parts) at a higher output capacitance per volt. Since the conduction loss dominates in this application (Section~\ref{sec:npar}), the per-volt view confirms the ranking. The price is a larger number of submodules (Fig.~\ref{fig:devices}c), i.e., more drivers and auxiliary supplies, which the modularity of the SPB accommodates~\cite{jin2017}. EPC2361 (100\,V, 0.75\,m$\Omega$ typical, 3\,$\times$\,5\,mm, top-side cooled) is selected. Its Kelvin source pad next to the gate returns the gate loop at the die, so that the inductance of the source stripes and vias stays outside the gate loop, a property that the gate drive of Section~\ref{sec:gate} relies on.

\begin{figure*}[t]
\centering
\includegraphics[width=\textwidth]{dev_compare.pdf}
\caption{Voltage-class selection for a 1200\,V SPB. (a) On-resistance against time-related output capacitance $C_\mathrm{oss,tr}=Q_\mathrm{oss}/V_\mathrm{test}$ (datasheet values; circles EPC, squares Infineon) with lines of constant $\tau=R_\mathrm{DS(on)}C_\mathrm{oss,tr}$. (b) Both quantities per volt of blocking capability. (c) SPBs from the best device of the 100, 200 and 650\,V classes, each axis normalized to its maximum: minimum loss~\eqref{eq:pmin}, devices per switch $n^\ast$, number of submodules, voltage step at the winding and number of transistors.}
\label{fig:devices}
\end{figure*}

\begin{table}[t]
\centering
\caption{Off-the-Shelf GaN Transistors for a 1200\,V SPB}
\label{tab:devices}
\setlength{\tabcolsep}{2.5pt}
\resizebox{\columnwidth}{!}{%
\begin{tabular}{@{}lccccccc@{}}
\toprule
Device & $V_\mathrm{DS}$ & $R_\mathrm{DS(on)}$ & $C_\mathrm{oss,tr}$ & $V_m$ & $\tau$ & $N$ & $P_\mathrm{min}$ \\
 & [V] & [m$\Omega$] & [pF] & [V] & [ps] & & rel. \\
\midrule
\DeviceRows
\bottomrule
\end{tabular}}
\end{table}

\subsection{Submodule voltage from the overshoot}
The submodule voltage plus the overshoot at turn-off must stay below 90\,V,
\begin{equation}
V_m+L_\mathrm{loop}\,\frac{di}{dt}\leq 90\,\mathrm{V}.
\label{eq:vm}
\end{equation}
Before a layout exists, $L_\mathrm{loop}$ can only be bounded. A vertical commutation loop over a 0.1\,mm dielectric with a few local MLCCs reaches about 0.4--0.5\,nH per cell, whereas a measurement shunt or a less compact layout gives about 1\,nH; both values are carried through the design. The turn-on $di/dt$, which the analytical gate-charge model of the 2EDF7275K-driven EPC2361 gives as \DidtSixtyA, \DidtEightyA{} and \DidtOneTwentyA\,A/ns at 60, 80 and 120\,A per device ($R_{g,\mathrm{on}}$\,=\,2\,$\Omega$), turns~\eqref{eq:vm} into the admissible submodule voltage of Fig.~\ref{fig:vm}(a). At 80\,A per device, which corresponds to four devices per switch at the peak current, \eqref{eq:vm} allows \VmaxEightyA\,V with 0.5\,nH but only \VmaxEightyB\,V with 1\,nH. With an even number of submodules, required for symmetric winding sets, $N$\,=\,16 and $V_m$\,=\,75\,V are selected (Fig.~\ref{fig:vm}b), and the layout must keep the loop near 0.5\,nH or the turn-on must be slowed. $N$\,=\,14 (86\,V) leaves no margin; $N$\,=\,18 (67\,V) costs 12.5\,\% more transistors, drivers and conduction loss, since every submodule carries the full phase current.

\begin{figure}[t]
\centering
\includegraphics[width=\columnwidth]{flow_vm2.pdf}
\caption{(a) Admissible submodule voltage $V_m=90\,\mathrm{V}-L_\mathrm{loop}\,di/dt$ for 0.5 and 1\,nH; the markers give the turn-on $di/dt$ at 60, 80 and 120\,A per device. (b) Submodule voltage $1200\,\mathrm{V}/N$; green: allowed and even.}
\label{fig:vm}
\end{figure}

%=========================================================================
\section{Number of Devices per Switch}\label{sec:npar}
%=========================================================================
\subsection{One device per switch}
A lumped double-pulse model of one EPC2361 per switch (LTspice, 75\,V, 1.7\,nH gate loop, 2\,$\Omega$ on, 0\,$\Omega$ off, 4/8\,A driver) bounds the current per device (Fig.~\ref{fig:single}). The peak $V_\mathrm{DS}$ stays at about 83\,V up to 80\,A with 0.5\,nH, but reaches 102\,V at 120\,A, and with 1\,nH already 93\,V at 80\,A (Fig.~\ref{fig:single}a). The turn-on $di/dt$ is set by the driver at about 45\,A/ns, whereas the turn-off $di/dt$ grows with the current (Fig.~\ref{fig:single}c). A larger loop lowers $E_\mathrm{on}$, because the loop takes $L_\mathrm{loop}\,di/dt$ of the voltage while the current rises and reduces the voltage--current overlap (28.3 versus 21.8\,\textmu J at 80\,A for 0.5 and 1\,nH), and it raises $E_\mathrm{off}$ and the overshoot; the energy $\frac12L_\mathrm{loop}I^2$ is dissipated in the ringing rather than in the transistor. The conduction loss per device is \PcondSixty, \PcondEighty{} and \PcondOneTwenty\,W at 60, 80 and 120\,A$_\mathrm{rms}$ and 100\,$^\circ$C (Fig.~\ref{fig:single}d). With \Ipk\,A peak per leg, a single device per switch is excluded, and the current per device is limited to about 80\,A.

\begin{figure*}[t]
\centering
\includegraphics[width=\textwidth]{flow_single2.pdf}
\caption{One EPC2361 per switch, lumped loop of 0.5 and 1\,nH (LTspice, 75\,V). (a) Peak drain--source voltage at turn-off; (b) turn-on and turn-off energy; (c) $|di/dt|$; (d) conduction loss per device against RMS current at 25, 100 and 150\,$^\circ$C.}
\label{fig:single}
\end{figure*}

\subsection{Devices in parallel and thermal path}
The number of devices per switch $n$ trades the conduction loss, which falls with $1/n$, against the switching loss, of which the output-charge part grows with $n$, and against the number of drivers. The loss of one leg is evaluated over the line period with the LTspice energies, with $T_j$ and $R_\mathrm{DS(on)}$ solved together for the thermal path of the load-test hardware: a Fischer LAM\,4\,K\,50\,12 fan heatsink (0.65\,K/W) on an aluminium spreader, and per device the junction-to-case resistance (0.2\,K/W), an insulating gap pad over the package top (\Rpad\,K/W) and the spreader path from the outer cells to the heatsink (\Rspr\,K/W), i.e., \Rjh\,K/W from junction to heatsink. Fig.~\ref{fig:npar} and Table~\ref{tab:npar} show the result at 50\,kHz. One device per switch runs away thermally, two reach \TjTwo\,$^\circ$C, and three keep $T_j$ at \TjThree\,$^\circ$C but exceed 90\,V with 0.5\,nH (\VpkThreeA\,V). Four devices per switch dissipate \TotFourA\,W per leg, i.e., a semiconductor efficiency of \EtaFour\,\% at 6.25\,kW per leg, at \TjFour\,$^\circ$C and 40\,$^\circ$C ambient, with \VpkFourA\,V peak voltage at 0.5\,nH and \VpkFourB\,V at 1\,nH. More devices gain less than 3\,W per leg but add drivers. Four devices per switch also split into two equal pairs, which the gate drive of Section~\ref{sec:gate} requires, and are therefore selected; the overshoot at 1\,nH confirms the layout target of 0.5\,nH. Conduction accounts for \CondFour\,W of the leg loss at 0.5\,nH and switching for \SwFourA\,W (\SwFourB\,W at 1\,nH); at 50\,kHz the leg is therefore conduction dominated, and the efficiency of the semiconductors falls from \EtaFTwenty\,\% at 20\,kHz to \EtaFFifty{} and \EtaFHundred\,\% at 50 and 100\,kHz.

\begin{figure*}[t]
\centering
\includegraphics[width=\textwidth]{flow_npar2.pdf}
\caption{$n$ EPC2361 per switch at \Irms\,A$_\mathrm{rms}$ and 50\,kHz. (a) Semiconductor loss per leg, left bars 0.5\,nH, right (hatched) 1\,nH; (b) junction temperature on the load-test heatsink for 25 and 40\,$^\circ$C ambient and a 3\,W/mK gap pad; (c) peak $V_\mathrm{DS}$ at the peak current per device.}
\label{fig:npar}
\end{figure*}

\begin{table}[t]
\centering
\caption{Semiconductor Loss per Leg (50\,kHz, 40\,$^\circ$C Ambient)}
\label{tab:npar}
\setlength{\tabcolsep}{3.5pt}
\begin{tabular}{@{}lccccc@{}}
\toprule
$n$ per switch & 2 & 3 & 4 & 6 & 8 \\
\midrule
Loss, 0.5\,nH [W] & \TotTwoA & \TotThreeA & \TotFourA & \TotSixA & \TotEightA \\
Loss, 1\,nH [W] & \TotTwoB & \TotThreeB & \TotFourB & \TotSixB & \TotEightB \\
Loss per device [W] & \PdevTwo & \PdevThree & \PdevFour & \PdevSix & \PdevEight \\
$T_j$ [$^\circ$C] & \TjTwo & \TjThree & \TjFour & \TjSix & \TjEight \\
$\hat v_\mathrm{DS}$, 0.5 / 1\,nH [V] & \VpkTwoA/\VpkTwoB & \VpkThreeA/\VpkThreeB & \VpkFourA/\VpkFourB & \VpkSixA/\VpkSixB & \VpkEightA/\VpkEightB \\
Driver channels per leg & 4 & 6 & 8 & 12 & 16 \\
\bottomrule
\end{tabular}
\end{table}

%=========================================================================
\section{Gate Drive and Current Sharing}\label{sec:gate}
%=========================================================================
Paralleled GaN transistors share the conduction loss well, because the on-resistance has a positive temperature coefficient, but not necessarily the switching loss. An LTspice double-pulse model of two paralleled EPC2361 (75\,V, 120\,A), in which one asymmetry is introduced at a time, quantifies the sensitivities (Fig.~\ref{fig:sharing}, Table~\ref{tab:sharing}). A gate-loop mismatch of 2 versus 8\,nH delays the second gate by about half of the current rise time and gives a dynamic imbalance of 19\,\% with a 4:1 split of the turn-on energy; without Kelvin connections a common-source inductance mismatch of only 0.2\,nH gives 29\,\%. A drain-path mismatch of 0.3\,nH, in contrast, gives 2\,\%. The gate loops and the common-source paths must therefore be equal by construction, whereas the power-loop paths are less critical.

\begin{figure*}[t]
\centering
\includegraphics[width=0.92\textwidth]{fig_sharing.pdf}
\caption{Current sharing of two paralleled EPC2361 (LTspice, 75\,V, 120\,A). (a) Drain currents at turn-on for (a.i) a symmetric layout with Kelvin connections, (a.ii) a gate-loop mismatch of 2 versus 8\,nH and (a.iii) no Kelvin connection with a 0.2\,nH common-source mismatch. (b) Turn-on (solid) and turn-off (hatched) energy of both transistors with the dynamic imbalance.}
\label{fig:sharing}
\end{figure*}

\begin{table}[t]
\centering
\caption{Sensitivity of the Switching-Energy Sharing (Two EPC2361)}
\label{tab:sharing}
\setlength{\tabcolsep}{3pt}
\resizebox{\columnwidth}{!}{%
\begin{tabular}{@{}llll@{}}
\toprule
Asymmetry & Case & Imbalance & Layout measure \\
\midrule
Power-loop inductance & $L_D$ +0.3\,nH & 2\,\% & mirrored cells \\
Gate-loop inductance & $L_G$ 2 vs.\ 8\,nH & 19\,\% & star gate network \\
Common source, no Kelvin & symmetric & 0\,\%, +47\,\% $E_\mathrm{on}$ & Kelvin pad \\
Common source, no Kelvin & $L_{CS}$ +0.2\,nH & 29\,\% & Kelvin return per device \\
Threshold spread & $V_\mathrm{th}$ +0.3\,V & 7\,\% & one channel per pair \\
\bottomrule
\end{tabular}}
\end{table}

A single driver channel for all four transistors of a switch is not a solution either. The four gate capacitances (4\,$\times$\,3.6\,nF) then charge through one driver output and a longer shared gate tree, so that $v_\mathrm{GS}$ rises four times more slowly, $di/dt$ and $dv/dt$ at turn-on fall and the overlap lasts longer: the turn-on energy per transistor grows by 16, 31 and 60\,\% for three, four and six transistors per channel, the peak sink current of 12--14\,A exceeds the 8\,A rating, and the longer tree increases the gate bump. The building block therefore uses one dual-channel isolated driver per pair of transistors, placed on the symmetry line between two mirrored cells and connected to both gates through a star with Kelvin returns (Fig.~\ref{fig:spb}b). Four devices per switch are built as two such pairs with two drivers per leg, i.e., 96 drivers for the complete inverter.

The 2EDF7275K (4/8\,A, 37\,ns propagation delay with at most 3\,ns channel-to-channel mismatch, 150\,V/ns common-mode transient immunity, 5\,$\times$\,5\,mm LGA) provides the high- and low-side channel of one pair in one package with isolated outputs. A 200\,V half-bridge driver with bootstrap supply (LMG1210, 1.5/3\,A class) is too weak for a pair, and single-channel isolated GaN drivers (1EDF5673K) double the number of driver packages. Each isolated output is supplied through a common-mode choke and split by a 5.1\,V Zener diode into $+5.1$\,V and a negative off-state voltage of about $-1.2$\,V. The negative bias is needed because the dv/dt of the complementary switching event lifts the gate of the off device by the Miller charge divided onto its gate--source capacitance, $\Delta v_\mathrm{GS}\approx Q_\mathrm{GD}/C_\mathrm{GS}\approx1.06$\,V, almost independent of current and bus voltage and close to the minimum threshold of 0.8\,V. With $-1.2$\,V the peak stays near 0\,V, the undershoot at $-2.3$\,V remains inside the $-4$\,V rating, and the larger reverse-conduction voltage in the dead time costs about 0.1\,W per leg.

%=========================================================================
\section{DC Link of a Three-Phase Level}\label{sec:dclink}
%=========================================================================
The three legs of a submodule share one dc bus. Because the phase currents of a balanced three-phase load sum to zero, the bus carries no ripple at twice the fundamental, and the capacitors only have to supply the switching-frequency component. A switching-function model of the three legs with space-vector modulation ($M$\,=\,1.15, $\cos\varphi$\,=\,0.9, \Ipk\,A peak) gives a capacitor RMS current of 85\,A per level, about \IripLeg\,A per leg, and a worst charge excursion within one switching period of 860\,\textmu C at 50\,kHz, close to $\hat I/(8f_\mathrm{sw})$. For a peak-to-peak ripple of 3\,\% (2.25\,V) the level needs 382\,\textmu F effective capacitance, i.e., 127\,\textmu F per leg board; the voltage ripple, not the RMS current, therefore sizes the dc link.

\begin{figure}[t]
\centering
\includegraphics[width=\columnwidth]{dclink2.pdf}
\caption{Number of 22\,\uF{} 2220 MLCCs per leg board for a given peak-to-peak ripple of the 75\,V level (three legs on one bus), with the dc-bias derating at 75\,V; star: selected design.}
\label{fig:dclink}
\end{figure}

Class-II MLCCs lose most of their capacitance at dc bias. The characteristic data of the manufacturers give only 15--29\,\% of the nominal value at 75\,V for 100\,V X7R/X7S parts of the 1210 to 2220 case sizes, well below the 30--45\,\% often assumed. The TDK C5750X7S2A226M280KB (2220, 22\,\uF) retains \CtwoRet\,\%, i.e., \CtwoEff\,\uF, the highest capacitance per part and per board area of the candidates, with \CtwoEsr\,m$\Omega$ ESR at 50\,kHz and \CtwoI\,A for a 10\,K temperature rise. Twenty parts per leg board in two rows meet the 127\,\textmu F (Fig.~\ref{fig:dclink}), each carrying about 1.4\,A of ripple current; the ripple current alone would need only \Nicap{} parts. Ten 1\,\uF{} 0805 capacitors per cell, six on the top side at the high-side drains and four directly below on the bottom side, supply the commutation edge. Their number follows from the loop inductance (beyond about four parts the ESL hardly matters, Section~\ref{sec:board}), and their capacitance from the hand-over of the load-current step to the bank through the inductance $L_b$ between the two, $C_\mathrm{byp}\geq L_b(\Delta I/\Delta V)^2$: with $L_b$\,=\,\LbMin--\LbMax\,nH from the Q3D model of the bank region and only 18\,\% of the nominal capacitance retained by the 0805 parts at 75\,V, ten parts per cell limit the dip to about 1.7\,V.

%=========================================================================
\section{Leg Board in Two Assemblies}\label{sec:board}
%=========================================================================
\subsection{Layout}
The leg board (89.5\,$\times$\,80\,mm, four layers, 2\,oz copper, 0.1\,mm dielectric between L1 and L2) carries the four cells of Fig.~\ref{fig:spb}(b) in two mirrored pairs with a driver in the centre of each pair (Fig.~\ref{fig:board}). Each commutation loop is vertical: the forward path on L1 closes on the dc$-$ plane on L2 directly below. One dc terminal pair (press-in M4 terminals on the bottom side) feeds a laminated terminal row (L1 and L3 dc$+$, L2 and L4 dc$-$) across the full board width, so that the bank of twenty 2220 capacitors and all four cells are supplied through the same copper structure. The switch nodes of the four cells are joined in copper on all layers and taken out at a single M6 bolt pad, to which a copper bar is bolted through a Hall-effect current transducer. A board generator creates the copper, the placement, the routing of the auxiliary circuits, the schematic and the fabrication data; the schematic is checked against the routed board net by net.

Two details proved essential for the symmetry of the cells. First, the vias of the source stripes cannot be placed in a row under the EPC2361, because the interleaved drain and source stripes leave no room between them; the source stripes of the low-side devices therefore carry vias in the pad. Second, every via that has to be removed for clearance must be removed together with its mirror image, otherwise the two cells of a pair end up with different return paths; with mirror-pair pruning the four loops agree within \LoopLoadSpread\,\% (Table~\ref{tab:q3d}).

\begin{figure*}[t]
\centering
\subfloat[]{\includegraphics[height=4.1cm]{r3d_dpt_top.png}}\hfill
\subfloat[]{\includegraphics[height=4.1cm]{r3d_load.png}}\hfill
\subfloat[]{\includegraphics[height=4.1cm]{q3d_loop_dpt.png}}
\caption{Leg board with four EPC2361 per switch (89.5\,$\times$\,80\,mm). (a) Double-pulse-test board with a coaxial shunt over every low-side transistor (top view with silkscreen). (b) Load-test board without shunts, with the 50\,mm fan heatsink on an aluminium spreader. (c) Q3D model of the cell band of the DPT board, exported from the routed copper (red dc$+$, blue dc$-$, orange switch node, green low-side source to the shunt; $z$ stretched).}
\label{fig:board}
\end{figure*}

\subsection{Two assemblies}
Measuring the current of every cell requires a shunt in each low-side source, which inevitably lengthens the commutation loop, whereas continuous operation needs a heatsink over the transistors, which the shunt bodies obstruct. The same design is therefore built in two assemblies. The double-pulse-test (DPT) board carries a coaxial shunt (PMK UFCS-R001, 1\,m$\Omega$, 120\,pH) per cell, whose blade is soldered to a land shortened to 1.2\,mm so that the blade sits 1.45\,mm from the source copper; it needs no heatsink. The load-test board omits the shunts and connects each low-side source through blind vias in the pad directly to the dc$-$ plane; it carries the heatsink of Section~\ref{sec:npar} and is the board used in the three-phase level. U.FL connectors on the bottom side give access to $v_\mathrm{GS}$ of both transistors and $v_\mathrm{DS}$ of the low-side transistor of every cell on both boards.

\subsection{Q3D verification}
To keep the extraction tractable, the routed board is split into small Q3D models, each solved only for the quantity it is used for: the cell band with the bypass capacitors and the shunts (ac at 100\,MHz) for the commutation loop, the full board (dc resistance) for the static current sharing and the copper loss, the cells with the joined ac copper for the inductance between the two driver pairs, and the terminal and bank region for the hand-over inductance $L_b$. The commutation loop is evaluated with all four low-side transistors switching, the high-side transistor of each cell shorted and the capacitor groups represented by their ESL.

\begin{figure}[t]
\centering
\includegraphics[width=\columnwidth]{loop_breakdown.pdf}
\caption{Commutation loop per cell of the final boards (Q3D). (a) Loop against the number of top-side bypass capacitors per cell (four more on the bottom side); stars: as built. (b) Contributions of the copper, the bypass ESL and the shunt, mean of the four cells.}
\label{fig:loop}
\end{figure}

\begin{table}[t]
\centering
\caption{Q3D Results of the Final Boards and the Busbar}
\label{tab:q3d}
\setlength{\tabcolsep}{3pt}
\resizebox{\columnwidth}{!}{%
\begin{tabular}{@{}lcc@{}}
\toprule
 & DPT board & Load-test board \\
\midrule
Loop per cell 1 / 2 / 3 / 4 [nH] & \LoopDptCOne/\LoopDptCTwo/\LoopDptCThree/\LoopDptCFour & \LoopLoadCOne/\LoopLoadCTwo/\LoopLoadCThree/\LoopLoadCFour \\
Copper / bypass ESL / shunt [nH] & \CuDpt/\EslDpt/\ShDpt & \CuLoad/\EslLoad/-- \\
$\hat v_\mathrm{DS}$ at 30\,A/ns [V] & \VosDpt & \VosLoad \\
Static sharing, max/min$-$1 (HS/LS) [\%] & \ImbDptHS/\ImbDptLS & \ImbLoadHS/\ImbLoadLS \\
Board copper loss at \Irms\,A$_\mathrm{rms}$ [W] & \PcuDpt & \PcuLoad \\
$L_x$ between the driver pairs (HS/LS) [nH] & \LxDptHS/\LxDptLS & \LxLoadHS/\LxLoadLS \\
\midrule
Busbar leg to leg, a--b / a--c, 1\,MHz [nH] & \multicolumn{2}{c}{\BusLab/\BusLac} \\
\bottomrule
\end{tabular}}
\end{table}

The loop is \LoopLoad\,nH per cell on the load-test board and \LoopDpt\,nH on the DPT board (Fig.~\ref{fig:loop}, Table~\ref{tab:q3d}). On the load-test board \CuLoad\,nH are copper and \EslLoad\,nH bypass ESL, so that the board meets the 0.5\,nH target with margin and gives about \VosLoad\,V peak voltage at 30\,A/ns. On the DPT board the copper path to the shunt adds 0.12\,nH and the shunt another 0.12\,nH; with \LoopDpt\,nH the turn-on must be limited to about \DidtDptMax\,A/ns for 90\,V, which is acceptable for a characterization board but confirms that the shunt does not belong in the continuously operated hardware. Beyond about four bypass capacitors per cell the loop no longer decreases noticeably (Fig.~\ref{fig:loop}a), and the bottom-side row lowers it by only 0.02\,nH; its purpose is capacitance, not inductance. The static sharing, i.e., the division of the conduction current among the four parallel cells set by their copper paths and $R_\mathrm{DS(on)}$, and the board copper loss follow from the dc model; the inductance $L_x$ between the switch nodes of the two driver pairs determines how fast a skew between the two drivers turns into circulating current, $\Delta I\approx V_m\Delta t/L_x$, and is the reason for driving both pairs from one PWM source.

%=========================================================================
\section{Three-Phase Level and Stack}\label{sec:level}
%=========================================================================
The three load-test boards of a level stand on their dc terminals (6.5\,mm high) on a laminated PCB busbar (298.5\,$\times$\,45\,mm, four layers, 2\,oz, L1 dc$+$ / L2 dc$-$ / L3 dc$+$ / L4 dc$-$ with 0.2\,mm prepreg), bolted from below with M4 screws (Fig.~\ref{fig:level}a). The busbar only carries the switching-frequency ripple between the banks of the three legs, since each commutation edge closes on its own board. For the banks to share this ripple, the inductance between neighbouring legs must be small against the capacitive reactance of a bank at 50\,kHz; the target of 5\,nH is met with \BusLab\,nH between neighbouring legs and \BusLac\,nH between the outer legs at 1\,MHz (\BusLabDc\,nH at dc, where the current spreads over all layers). Two M6 lug pads at the centre leg connect the level to the levels above and below, so that each busbar segment carries the dc current of only one leg. Every level floats up to the full stack voltage against the frame and is mounted on insulating standoffs with 8\,mm copper keep-out.

The leg pitch of 94.5\,mm (5\,mm gap) also fixes the heatsink: a 100\,mm fan heatsink with its fan is longer than the board and would collide with the heatsink of the neighbouring leg, whereas the 50\,mm version fits inside the board outline. Sixteen levels, each a three-phase winding set, form the 1200\,V stack (Fig.~\ref{fig:level}b); the dc$+$ lug of each level is joined to the dc$-$ lug of the next. The complete assembly is generated as a STEP model from the board data for the mechanical design.

\begin{figure}[t]
\centering
\subfloat[]{\includegraphics[height=3.5cm]{r3d_level_iso.png}}\hfill
\subfloat[]{\includegraphics[height=3.5cm]{r3d_stack.png}}
\caption{(a) One 75\,V level: three load-test leg boards on the laminated PCB busbar, each ac bar through its current transducer. (b) Sixteen levels in series form the 1200\,V stack (illustrative arrangement).}
\label{fig:level}
\end{figure}

%=========================================================================
\section{Discussion}
%=========================================================================
\begin{table}[t]
\centering
\caption{Comparison With Paralleled-GaN Half Bridges}
\label{tab:sota}
\setlength{\tabcolsep}{2.5pt}
\resizebox{\columnwidth}{!}{%
\begin{tabular}{@{}lllll@{}}
\toprule
Work & Devices/switch & Rating / test & $L_\mathrm{loop}$ & Layout \\
\midrule
\cite{lu2017} & 4\,$\times$\,650\,V & 400\,V/240\,A DPT & 0.7\,nH & 6-layer PCB \\
\cite{tran2025} & 4\,$\times$\,100\,V & 48\,V/328\,A DPT & 1.14\,nH & IMS hybrid \\
\cite{cooke2026} & 8\,$\times$\,200\,V & 120\,V/155\,A, 100\,kHz & 0.34\,nH$^\ast$ & PCB, symm.\ gate \\
\cite{wattenberg2023} & 8\,$\times$\,100\,V & 48\,V/125\,A$_\mathrm{rms}$ & -- & butterfly PCB \\
\cite{zou2025} & 3\,$\times$\,650\,V & 400\,V, 100\,kW & 6.4\,nH & PCB + Cu blocks \\
This work & 4\,$\times$\,100\,V & 75\,V/\Irms\,A$_\mathrm{rms}$ (design) & \LoopLoad\,nH/cell & four mirrored cells \\
\bottomrule
\multicolumn{5}{@{}l}{\footnotesize $^\ast$analytical estimate}
\end{tabular}}
\end{table}

Compared with the paralleled-GaN half bridges of Table~\ref{tab:sota}, the proposed leg reaches \LoopLoad\,nH per cell with four cells in parallel, and its symmetric gate drive removes the dominant sharing mechanisms of Table~\ref{tab:sharing}. Three limits remain. First, the margin to the 90\,V limit is set by the loop and the turn-on speed; with a measurement shunt the turn-on has to be slowed, and a regulated submodule voltage would be needed if the battery voltage rises well above nominal. Second, the thermal path of the test hardware is dominated by the insulating gap pad over each package, which the final integration would replace by a thinner, ceramic-insulated path~\cite{tran2025}. Third, at \Irms\,A$_\mathrm{rms}$ the copper loss of the board (\PcuLoad\,W, Table~\ref{tab:q3d}) adds to the semiconductor loss and has to be included in the thermal design. The two assemblies are prepared for the next step: double-pulse tests at 75\,V with the per-cell currents measured by the shunts, followed by continuous back-to-back operation of a three-phase level on the busbar.

%=========================================================================
\section{Conclusion}
%=========================================================================
A design procedure for the building block of stacked polyphase bridge converters has been carried from the specification of a 1200\,V, 300\,kW traction inverter to buildable hardware. Because the phase current is independent of the number of submodules, the device time constant and the on-resistance per volt select 100\,V GaN, and the overshoot at a realistic loop inductance sets 16 submodules of 75\,V. Loss, junction temperature and peak voltage give four EPC2361 per switch, arranged as two mirrored pairs with one dual-channel isolated driver each, and the dc link of a three-phase level is sized by the voltage ripple with the measured dc-bias derating, which selects twenty 22\,\uF{} 2220 MLCCs per leg. Q3D extraction of the routed boards confirms \LoopLoad\,nH per cell on the load-test board, \LoopDpt\,nH with the measurement shunt, equal loops in the four cells, and \BusLab\,nH between legs on the PCB busbar.

\begin{thebibliography}{99}
\bibitem{norrga2013} S.~Norrga, L.~Jin, O.~Wallmark, A.~Mayer, and K.~Ilves, ``A novel inverter topology for compact EV and HEV drive systems,'' in \emph{Proc. IECON}, 2013, pp. 6590--6595.
\bibitem{jin2014} L.~Jin, S.~Norrga, O.~Wallmark, and M.~Nikouei Harnefors, ``Control and modulation of the stacked polyphase bridges inverter,'' in \emph{Proc. IEEE ECCE}, 2014.
\bibitem{jin2017} L.~Jin, S.~Norrga, O.~Wallmark, and N.~Apostolopoulos, ``Modulation and power losses of a stacked polyphase bridge converter,'' \emph{IEEE J. Emerg. Sel. Topics Power Electron.}, vol.~5, no.~1, pp. 409--418, Mar. 2017.
\bibitem{nikouie2017tpel} M.~Nikouie, O.~Wallmark, L.~Jin, L.~Harnefors, and H.-P.~Nee, ``DC-link stability analysis and controller design for the stacked polyphase bridges converter,'' \emph{IEEE Trans. Power Electron.}, vol.~32, no.~2, pp. 1666--1674, Feb. 2017.
\bibitem{nikouie2017} M.~Nikouie, H.~Zhang, O.~Wallmark, and H.-P.~Nee, ``A highly integrated electric drive system for tomorrow's EVs and HEVs,'' in \emph{Proc. IEEE SPEC}, 2017.
\bibitem{zhang2014} H.~Zhang, O.~Wallmark, M.~Leksell, S.~Norrga, M.~Nikouei Harnefors, and L.~Jin, ``Machine design considerations for an MHF/SPB-converter based electric drive,'' in \emph{Proc. IECON}, 2014.
\bibitem{bringezu2020} T.~Bringezu and J.~Biela, ``Comparison of optimized motor-inverter systems using a stacked polyphase bridge converter combined with a 3-, 6-, 9-, or 12-phase PMSM,'' in \emph{Proc. EPE ECCE Europe}, 2020.
\bibitem{bringezu2023} T.~Bringezu and J.~Biela, ``Control of a 9-phase PMSM with stacked polyphase bridge converter including dc source impedance,'' in \emph{Proc. EPE ECCE Europe}, 2023.
\bibitem{verkroost2024} L.~Verkroost, F.~De~Belie, L.~Faria da Rocha, P.~Keim Olsen, P.~Sergeant, and H.~Vansompel, ``Multi-agent control scheme for power distribution in a multi-three-phase motor drive fed by a stacked polyphase bridges converter,'' \emph{IEEE Trans. Energy Convers.}, vol.~39, no.~4, pp. 2568--2580, Dec. 2024.
\bibitem{brothers2019} J.~A.~Brothers and T.~Beechner, ``GaN module design recommendations based on the analysis of a commercial 3-phase GaN module,'' in \emph{Proc. IEEE ECCE}, 2019, pp. 4109--4116.
\bibitem{lu2017} J.~Lu and D.~Chen, ``Paralleling GaN E-HEMTs in 10\,kW--100\,kW systems,'' in \emph{Proc. IEEE APEC}, 2017, pp. 3049--3056.
\bibitem{tran2025} M.~T.~Tran \emph{et al.}, ``A high-performance GaN power module with parallel packaging for high-current and low-voltage traction inverter applications,'' \emph{IEEE J. Emerg. Sel. Topics Power Electron.}, vol.~13, no.~1, pp. 1188--1209, Feb. 2025.
\bibitem{musumeci2022} S.~Musumeci, V.~Barba, and M.~Palma, ``GaN-based low-voltage inverter for electric scooter drive system,'' in \emph{Proc. AEIT Int. Annu. Conf.}, 2022.
\bibitem{wattenberg2023} M.~Wattenberg, O.~G.~Lorenz, and J.~Sanchez, ``A multi-kilowatt low-profile GaN inverter for light electric vehicles and high-power tools,'' in \emph{Proc. IEEE APEC}, 2023, pp. 1443--1450.
\bibitem{cooke2026} M.~W.~J.~Cooke and D.~J.~Rogers, ``120\,V, 155\,A traction inverter using eight parallel GaN HEMTs,'' 2026.
\bibitem{zou2025} M.~Zou \emph{et al.}, ``Systematic efficiency-density co-optimization of 100\,kW GaN traction inverter: methodology and integration,'' \emph{IEEE Trans. Power Electron.}, vol.~40, no.~11, pp. 16281--16299, Nov. 2025.
\bibitem{rohner2024} G.~Rohner, J.~Huber, S.~Miri\'c, and J.~W.~Kolar, ``Comparative evaluation of three-phase three-level flying capacitor and stacked polyphase bridge GaN inverter systems for integrated motor drives,'' \emph{Electronics}, vol.~13, no.~7, Art.~no.~1259, 2024.
\end{thebibliography}

\end{document}
